%!TEX root = ../LLMSeminar.tex
% ──────────────────────────────────────────────────────────────
%  Section 11 — The Bitter Lesson, Local Models, and Energy
% ──────────────────────────────────────────────────────────────

\section{The Bitter Lesson, local models, and the energy question}%
\label{sec:bitter-lesson}

We close the lecture with three pieces of homework.  They are
easy to do, individually small, and together they form a
calibration check for the rest of the seminar series.

\subsection{Homework 1: read \emph{The Bitter Lesson}}

\begin{keyresult}[Sutton's Bitter Lesson, in one sentence]
  ``The biggest lesson that can be read from $70$ years of AI
  research is that general methods that leverage computation
  are ultimately the most effective, and by a large margin.''
  \quad ---R.~Sutton, 2019~\cite{sutton2019bitter}.
\end{keyresult}

The essay is a single page.  Read it, then read it four more
times.  You should come out of it not entirely happy.

\paragraph{The thesis, in expanded form.}
Across seven decades of artificial-intelligence research, the
recurring pattern has been the following:

\begin{enumerate}[(i)]
  \item A research community decides that general methods
    are not enough and that progress requires baking in
    \emph{human insight}---hand-engineered features, expert
    knowledge bases, formal grammars, structured priors.
  \item Such systems achieve respectable, sometimes
    impressive, results.
  \item A few years pass.  Computers get larger.
  \item Someone runs a more general method---more search,
    more learning, less prior structure---at much greater
    scale, and outperforms the hand-engineered system by a
    margin that closes the case.
\end{enumerate}

The pattern repeats: in chess (Deep Blue), in Go (AlphaGo),
in speech recognition, in computer vision, in protein
folding, and---above all---in language modelling.  The bitter
content is this: the things people built lovingly with
domain knowledge ended up looking, in retrospect, like
detours.

\begin{intuition}[Why ``bitter'']
  The lesson is bitter because researchers naturally enjoy
  encoding their domain knowledge into systems; that is what
  expertise feels like in motion.  The empirical record is
  that the systems that win are those that resist the urge
  and let scale do the work.  Sutton's piece is uncomfortable
  reading for anyone who has spent a career on hand-built
  features.
\end{intuition}

\begin{remark}
  ``The bitter lesson'' is not the same as ``human insight is
  worthless.''  Sutton is careful: the insight that is durable
  is at the level of \emph{architecture for general
  computation}---what to learn, what to search, how to compose
  the two.  What does not survive is encoded prior knowledge
  about the specific domain.  The boundary is subtle and
  contested; argue with the essay seriously.
\end{remark}

\subsection{Homework 2: install Ollama and run a local model}

The second homework is concrete and physical: get a model
onto your own machine.  Ollama~\cite{ollama} is the easiest
on-ramp at the time of writing.  Download it, install it,
pull a model that fits in your memory, and have a brief
conversation.

\begin{algorithmbox}[Local-model setup, concretely]
  \begin{enumerate}[(1)]
    \item Visit \url{https://ollama.com}; install for your OS.
    \item Pick a model that fits your machine.  Rough rule:
      a model in $4$-bit quantisation occupies about half its
      parameter count in gigabytes of RAM (or VRAM).
      \begin{itemize}
        \item $\le 8$\,GB RAM: a $1$--$2$\,B parameter model
          (\texttt{qwen3:1.5b}, small \texttt{gemma3} variants).
          Useful for very small text tasks; do not expect
          arithmetic or careful reasoning.
        \item $16$--$32$\,GB RAM: $7$--$13$\,B models, e.g.\
          \texttt{qwen3}, \texttt{gemma3}, \texttt{llama3.x}.
          Begin to be useful for summarisation.
        \item $24$\,GB VRAM (e.g.\ a single RTX~3090):
          $13$--$30$\,B models comfortably; the larger ones run
          but are slow.
        \item $40$--$80$\,GB VRAM (e.g.\ A100/H100):
          frontier-class open-source models like
          \texttt{deepseek-r1}, \texttt{qwen3:72b}, large
          mixture-of-experts variants.  These can drive a
          coding agent.
      \end{itemize}
    \item Run \texttt{ollama run <model>} and interact.
    \item Try at least one summarisation task on a transcript
      or paragraph.
  \end{enumerate}
\end{algorithmbox}

\begin{remark}[Local-model competence profile]
  Local open-source models, even small ones, are reliably
  good at:
  \begin{itemize}
    \item Text summarisation.
    \item Cleaning up transcripts (especially Whisper output).
    \item Light editing and paraphrasing.
  \end{itemize}
  They are notably \emph{worse} than frontier proprietary
  models at:
  \begin{itemize}
    \item Arithmetic and quantitative reasoning---small
      models can second-guess themselves into nonsense on
      ``what is $1 + 1$?''
    \item Long-context coherence.
    \item Tool use, except for the specific models that have
      been trained for it.
  \end{itemize}
  The right use case for a local model is offline,
  text-heavy, low-stakes work---especially when the input is
  sensitive (interview transcripts, internal documents,
  patient discussions in medical settings) and must not
  leave the machine.
\end{remark}

\begin{audience}[Why text summarisation is a delicate test]
  Summarisation is unverifiable in the strict sense---to
  check the summary you have to read the original, which was
  the work the summary was supposed to save.  Nonetheless it
  is a legitimate use because the asymmetry favours us:
  \emph{producing} the summary takes longer than \emph{reading}
  it, even if checking the summary takes the same time as
  reading the source.  In practice, working from a summary
  with occasional reference to the source is faster than
  reading the source from cold.
\end{audience}

\subsection{Homework 3 (volunteer): the energy question}

Almost no first-order reasoning about LLMs in society proceeds
without someone asking what the things \emph{cost}, in energy
and carbon.  The popular numbers vary by orders of magnitude,
and few of them are well-sourced.  We need a careful look.

\begin{keyresult}[Volunteer task]
  Within a few weeks, a volunteer will produce a short
  presentation (no slides required) covering:
  \begin{enumerate}[(a)]
    \item \textbf{Inference cost.}  Measure, on your own
      laptop, the energy consumed by a typical query to a
      local model.  Express in joules per token.
    \item \textbf{Training cost.}  Find published figures
      for the training of a frontier model and convert to a
      reasonable amortisation per query, given typical
      lifetime usage estimates.
    \item \textbf{Carbon and water.}  Combine with average
      grid carbon intensity and data-centre water usage.
    \item \textbf{Social costs.}  The labour conditions of
      data labellers, the provenance of training data, and
      the dependence on unpaid (or under-paid) digital
      labour.
  \end{enumerate}
\end{keyresult}

\begin{remark}[Why this is non-trivial]
  The honest answer involves more than the marginal cost
  of a forward pass.  The accounting needs:
  \begin{itemize}
    \item training-time energy, amortised over the model's
      deployment lifetime;
    \item embodied energy in the GPU hardware
      (manufacturing footprint of an A100/H100 is
      substantial);
    \item cooling overhead (the data-centre Power Usage
      Effectiveness, or PUE);
    \item the ongoing cost of supporting infrastructure.
  \end{itemize}
  Compare this with claims such as ``one hour of LLM use
  costs the energy of a packet of crisps''---which may be
  numerically defensible if one counts only inference, but
  is not a complete accounting.
\end{remark}

\begin{historical}[A note on European efforts]
  Europe is in the middle of building its own foundation
  models under projects of which Zulfi (and the broader
  EuroLLM line of work) is one example, with the explicit
  goal of training on data sources whose provenance and
  labour conditions are documented.  Whether the next
  generation of European models is competitive on
  capability remains to be seen.  Whether they will be
  cleaner on the social-cost axis appears genuinely
  plausible.
\end{historical}

\subsection{Looking forward}

In the next seminar we will discuss two things.  First, the
content of \emph{The Bitter Lesson}---we will return to it
critically, with the experiments above as data.  Second,
\emph{proof assistance with AI tools}: what coding agents,
combined with the formal language Lean~4
\cite{moura2021lean}, can do for mathematical proofs, what
they cannot, and what the working pattern looks like.

\begin{audience}[A closing data point]
  An audience member reported, at the end of the seminar, a
  small experiment of their own: lab instruments
  (oscilloscope and waveform generator) connected to a
  laptop, with a frontier coding agent asked to identify an
  unknown circuit.  The agent designed a sequence of
  excitation signals, drove the instruments, observed the
  response, and correctly identified the device-under-test
  as a passive filter.  No human hint was given.  The
  remarkable feature, perhaps, is not that the agent
  succeeded but the strangeness of watching it
  \emph{actuate the physical world}---instrument relays
  clicking, signals propagating---driven by a chain of
  reasoning that lives entirely in tokens.
\end{audience}

\begin{intuition}[The unifying theme]
  Across all the experiments we have discussed---transcribing
  notes, building messengers, hijacking course schedules,
  inferring rooms from sound, identifying circuits from
  signals---the pattern is the same.  An agent's reach is set
  by its tools.  Give it more tools, and it touches more of
  the world.  The function~$f$ has not changed; only what
  the harness lets the function do.
\end{intuition}
